\section{The walk keeps its own clock}
\label{sec:clkwalk}

We know what the board says: some days carry more variance than
others.  This section shows that the mathematics has a natural home
for that fact.  The fact to establish, first in a discrete
computation a first course can check and then as a named theorem, is
this: \emph{the law of the underlying at expiry depends on the days
only through their total variance --- the budget --- and not on how
that budget is scheduled.}  Once that is true, ``which clock to use''
becomes a genuine choice, because the prices cannot object.

Do the discrete case in full.  Let the log return of the underlying
over trading day $i$ be a random shock $\varepsilon_i$, normally
distributed with mean zero and variance $\Delta w_i$ --- day $i$'s
contribution to the total.  (Mean zero is a deliberate
simplification: it drops the $-\tfrac12$ drift correction Chapter~5
derived, which shifts the mean and not the variance; the correction
reappears, intact, in the continuous computation below.)  Let the
days be independent: today's surprise does not read yesterday's.  After $n$ days the cumulative log return is the sum
$\varepsilon_1+\cdots+\varepsilon_n$, and two first-course facts
finish the argument.  A sum of independent normals is normal.  And
its variance is the sum of the variances:
\[
\operatorname{Var}\!\big(\varepsilon_1+\cdots+\varepsilon_n\big)
=\Delta w_1+\Delta w_2+\cdots+\Delta w_n .
\]
The right-hand side is a \emph{sum}, and sums do not care about order
or arrangement.  Put the big day first, last, or in the middle; split
it into two medium days; the terminal law is the same normal with the
same variance.  The law at expiry knows the budget, not the schedule.

Make it concrete with the numbers this section's figure uses:
\MacClkWalkDays\ trading days, each ordinary day's shock worth
$\MacClkWalkDailyPct\%$ of standard deviation --- that is,
$(\MacClkWalkDailyPct\%)^2$ of variance --- and one scheduled event on
day \MacClkWalkEventDay\ carrying \MacClkWalkEventUnits\ ordinary
days' variance in one calendar day.  Take one ordinary day's variance
as the unit.  The \MacClkWalkDays\ days are 29 ordinary ones plus the
event day; the ordinary days contribute one unit each and the event
day contributes \MacClkWalkEventUnits, so the terminal standard
deviation is
\[
\MacClkWalkDailyPct\%\times
\sqrt{\underbrace{29\times1}_{\text{ordinary days}}
+\underbrace{1\times\MacClkWalkEventUnits}_{\text{the event day}}}
=\MacClkWalkDailyPct\%\times\sqrt{34},
\]
and it would be $\MacClkWalkDailyPct\%\times\sqrt{34}$ with the event
on day 3, day 20, or day 30.  An option struck on the terminal level
--- any option, since Chapter~2 taught that a full strike ladder pins
the terminal law and nothing more --- cannot tell those calendars
apart.

But something \emph{can} tell them apart: any reading that divides by
elapsed calendar time.  \Cref{fig:clkwalk} shows the same simulated
walk twice.  In panel~(a), the horizontal axis is the trading day.
The blue band with grey edges is the $\pm1$-standard-deviation
envelope, $\pm\MacClkWalkDailyPct\%\sqrt{\text{variance so far}}$:
the statistical corridor the walk lives in.  Follow the band's upper edge
with your eye: it grows like a square root, then takes a visible
step up at the dotted line --- day \MacClkWalkEventDay, where one
day adds five days' variance --- and resumes its square-root growth
from a higher base.  The kink is the calendar clock failing to keep
pace with the walk.  In panel~(b), the \emph{same} paths are drawn
against a different horizontal axis: accumulated variance, measured
in \emph{variance days} --- ordinary days count one, the event day
counts \MacClkWalkEventUnits.  The event day is now simply a wider
slot on the axis (the orange strip), the envelope is one smooth
square root from end to end, and nothing about the walk marks the
event at all.  Same randomness, same terminal law; the kink lived in
the ruler.

\begin{figure}[htbp]
\centering
\paperfig{\textwidth}{fig_clk_walk.pdf}
\caption{One walk, two rulers (\MacClkWalkPaths\ simulated paths, one
highlighted; construction and seed in
appendix~\ref{app:clkprotocol}).  (a)~Against the trading-day index,
the $\pm1$-sd envelope (grey edges, blue band) kinks at the event
day: day \MacClkWalkEventDay\ carries \MacClkWalkEventUnits\ days'
variance.  (b)~Against accumulated variance the event day is a wider
slot on the axis (the orange strip), and the envelope is one smooth
square root.  The walk itself never noticed the event's position;
only the ruler did.}
\label{fig:clkwalk}
\end{figure}

The continuous version of this additivity is a classical theorem.
Chapter~5 wrote the underlying's motion as a diffusion,
$\dd Y_s=\sqrt{v_s}\,Y_s\,\dd W_s$ --- $Y$ the forward-normalized
level, $W$ a Brownian motion (the continuous limit of an unbiased
random walk), $v_s$ the instantaneous variance, the squared
volatility the process runs at instant $s$ (\cref{sec:vsnumber}).
There the time variable was variance time itself; run it now on the
\emph{calendar}, where $v_s$ is free to schedule itself unevenly ---
quiet Tuesdays, a violent earnings evening.  Chapter~5's It\^o
computation \eqref{eq:vsito} splits the log return into a drift and a
martingale part --- a \emph{martingale} being, as before, a process
with no drift, whose expected future value is its current value:
\[
\log Y_t
=\underbrace{\int_0^{t}\sqrt{v_s}\,\dd W_s}_{\textstyle M_t,
\ \text{the martingale part}}
\;-\;\frac12\underbrace{\int_0^{t}v_s\,\dd s}_{\textstyle
\mathcal{C}(t),\ \text{accumulated variance}} .
\]
The second integral, written $\mathcal{C}(t)$, is the continuous
analogue of ``variance days'': the running total of instantaneous
variance, growing fast through violent stretches and slowly through
quiet ones.  It is the clock the market itself keeps.  For a general
continuous martingale --- one not handed to us as an integral against
a Brownian motion --- the accumulated variance is defined directly
from the path, as the small-step limit of the sum of squared
increments; for our $M_t=\int_0^t\sqrt{v_s}\,\dd W_s$ that limit is
exactly the $\mathcal{C}(t)=\int_0^t v_s\,\dd s$ above.  With the
term defined, the theorem says the martingale part is nothing but a
Brownian motion read on that clock:

\begin{theorem}[Dambis; Dubins--Schwarz
\citep{Dambis1965,DubinsSchwarz1965}]
\label{thm:clkdds}
Let $M$ be a continuous martingale with $M_0=0$ whose accumulated
variance $\mathcal{C}$ grows without bound.  Then there is a standard
Brownian motion $Z$ such that
\[
M_t \;=\; Z_{\mathcal{C}(t)}\qquad\text{for all }t .
\]
\end{theorem}

Two remarks set the statement's scope at this book's level.  The
growth condition just ensures the relabeled clock runs on forever, so
that $Z$ is a full Brownian motion; it is harmless here, where the
market keeps moving past any one expiry.  And the delicate step ---
reading a Brownian motion at the \emph{random} time $\mathcal{C}(t)$
and getting the original martingale back --- is precisely what the
proof establishes, by a careful change-of-time argument in the
standard reference \citep[Ch.~V]{RevuzYor1999}.  We use the
statement, not the proof.
What it buys us is the discrete computation again, now for every
continuous martingale at once: running the motion to calendar date
$t$, the log return is
$\log Y_t=Z_{\mathcal{C}(t)}-\tfrac12\mathcal{C}(t)$, so when the
schedule $v_s$ is deterministic the terminal law is normal with
variance $\mathcal{C}(t)$ --- the budget --- and mean
$-\tfrac12\mathcal{C}(t)$, no matter how the accrual was arranged
along the way.  The market keeps one clock, $\mathcal{C}$; the
calendar is a second clock that we brought.

Now say what this means for implied volatility, in one careful
paragraph.  A quoted option price, pushed through the Black formula
of Chapter~2, delivers a \emph{total variance} $w(k)$ --- and that
inversion never asks about time at all.  This is worth pausing on,
because a first course writes Black--Scholes with $\sigma$ and $T$ as
separate inputs.  Chapter~2's normalized form \eqref{eq:black} shows
they never act separately: the formula is $B(k,w)$ with
$d_\pm=-k/\sqrt{w}\pm\sqrt{w}/2$, and maturity enters only through
the product $w=\sigma^2\tau$.  Everything after the inversion is
division: quoting the same $w$ as a volatility requires choosing a
clock and computing $\sqrt{w/\text{clock}}$.  Divide by the calendar
$t$ and we write
\[
\sigma_{\rm cal}=\sqrt{w/t}\,;
\]
divide by a clock $\tau$ built to follow the market's own accrual
schedule and we get $\sigma=\sqrt{w/\tau}$ --- which is precisely the
$\sigma$ this book has used since Chapter~2, where $w=\sigma^2\tau$
was the founding convention.  Until today the two clocks agreed,
because no event calendar was in force; the NVDA board is where they
part.  The dip of \cref{sec:clkboard} is now diagnosed: the 18-day
expiry dodges a scheduled burst of $\mathcal{C}$, its $w$ is
accordingly modest, and dividing that modest $w$ by an
\emph{undiscounted} 18 days of calendar produces a number that looks
like a volatility collapse.  The prices are consistent; the reading
$\sigma_{\rm cal}$ is the artifact.

One boundary before building the clock.  The theorem licenses
a \emph{relabeling} of time; it does not predict $\mathcal{C}$.  The
market's clock is random --- nobody knows Tuesday's variance in
advance --- and what a fitting desk can carry is a deterministic
\emph{forecast of its schedule}: ordinary days advance the clock at
rate one, scheduled event days advance it by more.  Building that
forecast, as a small object a desk can quote and edit, is the next
section's job.
